\section{The on-chain solver}
\label{sec:solver}

\Cref{prop:E,prop:E-strict} establish that the fixed-point equation
\eqref{eq:fixed-point} has a unique root in the bracket
$[N\alpha,\,N\beta]$ whenever the loop gain $g$ is below $1$ and the
error function changes sign across the bracket. This section specifies how
the kernel finds that root on-chain: bisection on $E(x)$, with bounded
gas, deterministic iteration count, conservative rounding, and explicit
revert paths when any precondition fails.

\subsection{Integer bisection on $E(x)$}
\label{sec:solver:bisection}

The solve operates on $E$ at integer (WAD) precision. At each iteration
the bracket is halved and the half containing the sign change is kept.
The loop runs for a fixed number of steps determined by the bracket
width and the configured tolerance, then returns the midpoint (rounded
toward senior, see \S\ref{sec:solver:rounding}) as $x^*$. The published rate is
$y^* = x^* / N$.

Bisection is the right primitive here for three reasons:
\begin{enumerate}
  \item \textbf{No derivative needed.} The waterfall $F$ is non-smooth
    (\cref{lem:F}), so any tangent-following method (Newton, secant)
    can step across a kink incorrectly. Bisection requires only that
    $E$ be continuous and strictly monotone, both established in
    \S\ref{sec:existence}.
  \item \textbf{Deterministic gas.} The iteration count depends only on
    the bracket width and the tolerance, not on the input data, so
    every sync costs the same regardless of pool state. This makes the
    cost of an operation predictable from the caller's side.
  \item \textbf{Auditable.} The solver is a short loop over already-audited
    valuation code ($P_A$, the waterfall $F$). There is no closed-form
    invariant inversion to verify, no spurious-root selection logic, and
    nothing pool-specific in the solver itself.
\end{enumerate}

\subsection{Default bracket}
\label{sec:solver:bracket}

The default search bracket is the rate range the pool can express. With
$\alpha$ and $\beta$ the lower and upper bounds on the senior share
price in quote (the orientation used throughout this document), the
senior rate satisfies $y \in [\alpha, \beta]$, so
\begin{equation}
  x \in \left[N\alpha,\, N\beta\right].
  \label{eq:default-bracket}
\end{equation}
On-chain, the E-CLP's parameters bound the price of its token0 in units
of its token1, and Balancer orders tokens by address: if the deployed
pool sorts the senior share as token1, the configured E-CLP bounds are
the reciprocals of the $\alpha, \beta$ used here. Both bounds are
per-solve quantities. Because the rate is injected through the senior
leg's balance, the configured band constrains the rate-scaled price, so
the senior-price bounds used here are the configured band recentered at
the pre-solve cached rate.

This is the operating range of the pool, not a universal economic bound.
A senior loss large enough to push the correct rate below $\alpha$ falls
outside the band the pool can represent. The solver must reject such
states rather than clip the rate to the band edge
(\S\ref{sec:solver:failure}).

A production market may seed a tighter sub-bracket around an estimate of
the operating rate to keep the iteration count small, but the wider
$[N\alpha, N\beta]$ is the conservative default. Either way, the
bracket is fixed at the start of each solve from values read once, and
iterations do not re-read pool state.

\subsection{Sign-change requirement and failure modes}
\label{sec:solver:failure}

By \cref{prop:E-strict}, $E$ is strictly decreasing on the bracket. A
unique root exists iff $E$ has opposite signs at the endpoints:
\[
  E(x_{\mathrm{lo}}) \geq 0 \;\geq\; E(x_{\mathrm{hi}}).
\]
The solver checks this condition at the start of every solve. There are
four hard-revert paths the kernel takes when a precondition fails. In
each case the cached rate is left untouched, so a stale value cannot be
read after a failed solve:
\begin{description}
  \item[\code{NO_SIGN_CHANGE}] $E$ has the same sign at both bracket
    endpoints. The correct rate sits outside the band the pool can
    express, so the sync reverts.
  \item[\code{MAX_ITER_EXCEEDED}] The configured iteration count was
    insufficient to drive the residual below tolerance. Should not occur
    under the chosen bracket and tolerance, but reverts as a safety net.
  \item[\code{RESIDUAL_OUT_OF_TOLERANCE}] The final residual is above the
    configured ceiling. Same posture as above.
  \item[\code{LOOP_GAIN_EXCEEDED}] The loop-gain proxy
    $g_{\text{proxy}} = f \cdot \varphi$ is at or above
    \code{maxLoopGainWAD} at solve time
    (\S\ref{sec:solver:loop-gain}).
\end{description}

A governance-approved bounded fallback bracket may be configured for the
\code{NO_SIGN_CHANGE} case. If configured, the solver re-attempts within
the wider band. If the wider band also lacks a sign change, the sync
still reverts. The solver never returns a bracket endpoint and never
overwrites a previously-published rate after a failed solve.

\subsection{Tolerance, iteration count, and rounding}
\label{sec:solver:rounding}

Three policy choices fix the solver's deterministic behaviour:
\begin{description}
  \item[Tolerance] the residual ceiling at which the solver stops
    halving. Should be at or above the accountant's dust tolerance:
    resolving below dust changes nothing the accountant can observe.
  \item[Iteration count] fixed at
    $\lceil \log_2(\mathrm{width} / \mathrm{tolerance}) \rceil$.
    Hardcoded per market at deploy time, so gas is data-independent and
    the loop has no conditional break. Width is taken from
    \cref{eq:default-bracket} and tolerance from the kernel's storage.
    The width and residual readings of tolerance coincide: $E$ has
    slopes in $[-1,\ g-1]$ on the bracket, so it is 1-Lipschitz, and
    driving the bracket width below the tolerance drives the residual
    below it as well.
  \item[Rounding direction] the final $x^*$ is rounded toward the senior
    side. Together with floor-rounding inside $P_A$
    (\eqref{eq:PA}), this ensures any 1-wei accumulation favours the
    protected tranche.
\end{description}

\subsection{Loop-gain guard}
\label{sec:solver:loop-gain}

\Cref{prop:E-strict} requires $g < 1$ on the search bracket. The true $g
= f \cdot \sup_y \mathrm{yb}(y) / N$ requires evaluating $\mathrm{yb}$
across the band, which is more expensive than the rest of the sync.
The kernel substitutes a cheap proxy:
\[
  g_{\text{proxy}} \;=\; f \cdot \varphi,
    \qquad
  \varphi = b_s / N.
\]
At the balanced point $\mathrm{yb}(y) \approx b_s$, so $g_{\text{proxy}}
\approx g$. Away from the balanced point (particularly near the
low-rate end of a wide band) the true $g$ can exceed $g_{\text{proxy}}$
because the E-CLP geometry amplifies rate changes. $g_{\text{proxy}}$ is
therefore a \emph{lower-order risk indicator}, not an upper bound on the
worst-case loop gain.

Governance configures a ceiling \code{maxLoopGainWAD}. Before solving,
the kernel asserts $g_{\text{proxy}} <$ \code{maxLoopGainWAD} and
reverts the sync with \code{LOOP_GAIN_EXCEEDED} otherwise. Because the proxy
underestimates the worst-case gain, the production ceiling should sit
\emph{strictly} below WAD with sufficient margin, and may be paired
with hard caps on $f$ and $\varphi$, or with a stronger endpoint
evaluation of $g$ at the bracket ends. The precise margin is a Royco risk
parameter.

\subsection{Properties auditors should verify}
\label{sec:solver:audit-properties}

The solver implementation is required to satisfy:

\begin{enumerate}
  \item the iteration count is a deterministic function of bracket and
    tolerance, hardcoded per market, with no data-dependent early exit;
  \item the bracket is computed once per solve from values read at the
    start, and iterations do not re-read pool state (no live calls back
    into the Vault during the solve);
  \item every revert path
    (\code{NO_SIGN_CHANGE}, \code{MAX_ITER_EXCEEDED},
    \code{RESIDUAL_OUT_OF_TOLERANCE}, \code{LOOP_GAIN_EXCEEDED}) leaves
    the cached rate untouched;
  \item the final $x^*$ is rounded toward the senior side;
  \item the solver never returns a bracket endpoint as the solution.
\end{enumerate}
